\section{量词与否定}
\label{sec:01-Mathematical-Language/01-Logic/03}

有人扔给你一句话：``\(x^2\ge 0\)。''问你，这句话对不对？

你可能会说：当然对，平方不可能小于零。但等一下——如果 \(x\) 是复数呢？\(i^2=-1\)，小于零。对方没说过 \(x\) 是实数。你也没问。

这就是本节要处理的问题。上一节用真值表和蕴含箭头处理了命题之间的连接，但那些命题本身必须是完整的。``\(x^2\ge 0\)''根本不完整，它缺了一个关键零件：\textbf{\(x\) 从哪来的，以及你要对多少个 \(x\) 下断言}。

\subsection{一个不完整的句子，加上三件东西才成命题}

单独一句``\(x^2\ge 0\)''，上一节叫它开放命题或者谓词 \(P(x)\)。你没法判断它真假，因为 \(x\) 悬在半空。要让这句话着陆，需要三样东西：

\begin{itemize}
\tightlist
\item
  \textbf{变量}是谁——这里是 \(x\)；
\item
  \textbf{论域}——\(x\) 从哪个池子里取，是 \(\R\)？\(\Z\)？还是 \(\C\)？
\item
  \textbf{量词}——你要断言论域里的\textbf{每一个}都满足，还是\textbf{至少有一个}满足？
\end{itemize}

三样凑齐：

\[
\forall x\in\R,\quad x^2\ge 0.
\]

现在它可以判断真假了：对实数集里的每一个 \(x\)，它的平方都不小于 0。真的。

把论域换一下试试。同样一句话``\(x^2=2\)''，套上不同的论域，结果完全不同：

\begin{itemize}
\tightlist
\item
  在整数集 \(\Z\) 里：\(\exists x\in\Z,\,x^2=2\)——假的，没有一个整数的平方是 2。
\item
  在实数集 \(\R\) 里：\(\exists x\in\R,\,x^2=2\)——真的，\(\sqrt{2}\) 和 \(-\sqrt{2}\) 都是。
\item
  在正实数集 \(\R^+\) 里：\(\exists!x\in\R^+,\,x^2=2\)——唯一存在，只有 \(\sqrt{2}\)。
\end{itemize}

同一个谓词，池子不同，真假就不同。论域不是装饰，是命题的半个身子。写量词的时候心里同时装着三件事：变量是谁、池子在哪、断了什么性质。

\subsection{怎么证明，怎么推翻}

全称命题 \(\forall x\in D, P(x)\) 的要害在于：它要求池子里\textbf{每一个}都满足。所以想证明它，你得对池子里的任意一个成员做通行的论证——不能只查几个样本就宣布胜利。而推翻它就极其简单：找出一个反例就行。一个就够了。

反过来，存在命题 \(\exists x\in D, P(x)\) 只要找到一个例子就能证明；推翻它却是一场苦役——你得说明池子里的每一个成员都不满足 \(P(x)\)。

这一进一退，就是量词否定律：

\[
\neg\forall x,P(x)\equiv\exists x,\neg P(x),
\qquad
\neg\exists x,P(x)\equiv\forall x,\neg P(x).
\]

翻译回日常语言：``不是所有人都及格''等于``至少有一个人没及格''；``没有人及格''等于``每一个人都没及格''。听起来像废话，但在复杂嵌套里操作时这套规则是救命绳——否定一个带量词的命题，量词要翻转（\(\forall\leftrightarrow\exists\)），后面的性质也要否定，两步缺一不可。

拿上一节用过的东西练手：``不是每个连续函数都可微''。按规则，\(\neg\forall f\)（连续 \(\to\) 可微）\(\equiv\) \(\exists f\)（连续 \(\land\) 不可微）。\(f(x)=|x|\) 在 \(x=0\) 连续但不可微——它就是那个见证。注意，否定的结果是``存在连续不可微的函数''，不是``每个连续函数都不可微''。这两个差了一个量词，差了一整个世界。

\subsection{量词顺序：先拿门票还是先选座位}

这是量词最要害的地方，也是新手翻车最密集的路口。看这两句：

\[
\forall x\in\R,\; \exists y\in\R,\; y>x,
\]

和

\[
\exists y\in\R,\; \forall x\in\R,\; y>x.
\]

看起来只是把 \(\forall x\) 和 \(\exists y\) 换了个位置，意思却天差地别。

第一句：给你一个 \(x\)，比如 \(x=100\)，你就在它后面找个更大的 \(y\)，取 \(y=x+1=101\) 就行。换 \(x=10^6\)？取 \(y=10^6+1\)。每个 \(x\) 后面你都\textbf{允许重新选}一个 \(y\)。总是成立。

第二句：你得先找到一个固定的 \(y\)，然后这个 \(y\) 得比\textbf{所有}实数都大。你试试看：取 \(y=10^{100}\)？不够，\(x=10^{100}+1\) 比它大。取 \(y\) 任意大？总有比它大一号的实数。不存在这样的 \(y\)。命题是假的。

关键在于依赖关系的方向。\(\forall x\exists y\) 的意思是：\textbf{\(y\) 可以依赖 \(x\)}——每个 \(x\) 配一个专属的 \(y\)。\(\exists y\forall x\) 的意思是：\textbf{\(y\) 必须预先选定，然后独立面对所有的 \(x\)}——同一个 \(y\) 要能扛住整个池子。

拿一串钥匙就能看明白。设 \(L\) 是家里的锁（前门、信箱、自行车锁、地下室锁），\(K\) 是钥匙串上的钥匙。``每把锁都有钥匙能打开''：

\[
\forall \ell\in L,\; \exists k\in K,\quad k\text{ 能打开 }\ell.
\]

这当然成立——你拿起前门的锁，找对应的钥匙；拿起信箱的锁，找另一把钥匙。不同的锁用不同的钥匙，没毛病。

但如果交换量词：

\[
\exists k\in K,\; \forall \ell\in L,\quad k\text{ 能打开 }\ell.
\]

这就成了``存在一把万能钥匙，能打开家里所有的锁''。你可能立刻会说：哪有这种好事。对，通常不成立。两句话的区别，就是你愿不愿意掏一串钥匙去一个个试，和你有没有一把能通吃的钥匙，这两种生活经验的区别。

写证明的时候，把这种依赖关系亮出来。如果你证的是 \(\forall x\exists y\) 的形式，写清楚``给定 \(x\)，取 \(y=\underline{\quad}\)''，让读者看见 \(y\) 是怎么根据 \(x\) 构造出来的。别只甩一句``显然存在 \(y\)''，那等于什么都没说。

同类型量词交换顺序没事：\(\forall x\forall y\) 和 \(\forall y\forall x\) 说的都是``每一对 \((x,y)\) 都满足''，一个意思。\(\exists x\exists y\) 也一样，换顺序不改变存在一对见证这件事。但\textbf{不同类型}的量词，一般情况下不可交换。每当你看到 \(\forall\exists\) 或者 \(\exists\forall\)，脑子里先停半秒：谁选谁？谁先固定，谁后匹配？

\subsection{条件句里的量词和空真}

回到一个熟悉的句子：``每个偶数都是 4 的倍数。''你已经知道这是假的——2 是偶数但不是 4 的倍数。把它用量词写出来：

\[
\forall n\in\Z,\quad n\text{ 是偶数}\to n\text{ 是 4 的倍数}.
\]

这句话要求的是：对\textbf{每一个}整数 \(n\)，蕴含 \(P(n)\to Q(n)\) 都成立。检查 \(n=2\)：前提``2 是偶数''真，结论``2 是 4 的倍数''假。所以 2 是一个反例，全称命题被推翻。顺便说一句，奇数 \(n\) 不会带来麻烦——前提假了，蕴含自动真，这是上一节真值表就定好的事。

当然你也可以直接把论域抠成偶数集合再写全称，效果一样。两种写法等价，选你觉得顺手的。

但如果整个论域里\textbf{根本没有}满足前提的元素呢？比如：

\[
\forall x\in\Z,\quad x>10\land x<11\to x\text{ 是蓝色的}.
\]

论域是整数。存在大于 10 且小于 11 的整数吗？不存在。所以前提对每个整数都是假的——整句话因为在论域里找不到反例而自动成立。这叫\textbf{空真}。

空真听起来像逻辑学家的诡辩：``所有大于 10 的一位正整数都是蓝色的''居然算真命题？但这是形式系统运转的代价。有一个你每天在用的东西靠空真活着：\textbf{空集是任何集合的子集}。\(A\subseteq B\) 的定义是 \(\forall x(x\in A\to x\in B)\)。当 \(A=\varnothing\) 时，前提 \(x\in\varnothing\) 永远假，蕴含全真，所以 \(\varnothing\subseteq B\) 对所有 \(B\) 成立。要是不承认空真，这条最基本的性质就得打上一堆补丁，更麻烦。

空真给我们的提醒是：一个全称条件命题的真，不总是意味着你获得了关于世界的正面知识。有时候它只是在告诉你——\textbf{被量化的那批对象是空的}。

\subsection{唯一存在：存在不够，还得没有第二个}

\(\exists!x,P(x)\) 比 \(\exists x,P(x)\) 多了一个感叹号，也多了一倍的证明工作量。你不仅要找到一个满足 \(P(x)\) 的对象，还得论证\textbf{没有第二个}。

比如：对任意实数 \(a\)，方程 \(x+a=0\) 有唯一实数解。存在性：取 \(x=-a\)，代进去 \(-a+a=0\)，没问题。唯一性：假设 \(x_1\) 和 \(x_2\) 都是解，即 \(x_1+a=0\) 且 \(x_2+a=0\)。两式相减，\(x_1-x_2=0\)，所以 \(x_1=x_2\)。没有第二个候选。

标准的唯一性证明模板就是这两步：先造（或论证）一个，再假设有两个，推出它们相等。光写出一个解就收工，排除不了藏着另一个解的可能。

再看一个需要警觉的例子。``方程 \(x^2=4\) 在实数里有唯一解''——不对，因为 \(2\) 和 \(-2\) 都是解。但如果你把论域限定为``正实数''，那确实只有 \(2\)。再一次看到，\textbf{论域变了，断言的力量就变了}。量词的强弱，一半在量词本身的种类，一半在论域的大小。

\subsection{否定复杂命题：从外向内一层一层翻}

来看一个以后你会反复遇到的定义——一致连续。你现在不需要理解它的几何含义，只需要把它当成一条嵌套量词命题，练一次否定它的操作：

\[
\forall\varepsilon>0,\; \exists\delta>0,\;
\forall x,y\in D,\quad
|x-y|<\delta\to|f(x)-f(y)|<\varepsilon.
\]

否定嵌套量词命题，方法就是\textbf{从最外层开始，一层一层往里翻}。每翻一层，量词翻转（\(\forall\leftrightarrow\exists\)），然后继续处理里面那一层，直到最内层的谓词。每步只做一件事，不跳。

\begin{enumerate}
\tightlist
\item
  外层是 \(\forall\varepsilon>0\)。翻转：\(\exists\varepsilon_0>0\)。下标 \(_0\) 只是标记这是被挑出来的那个``坏''数，不是变量名有特殊含义。
\item
  下一层是 \(\exists\delta>0\)。注意这一层在 \(\exists\varepsilon_0\) 的\textbf{里面}，所以否定它：翻成 \(\forall\delta>0\)。
\item
  再下一层是 \(\forall x,y\in D\)。否定它：翻成 \(\exists x,y\in D\)。
\item
  现在到了最内层的蕴含：\(|x-y|<\delta\to|f(x)-f(y)|<\varepsilon\)。上节讲过的，\(\neg(P\to Q)\equiv P\land\neg Q\)。所以翻成：
  \[
  |x-y|<\delta\;\land\;|f(x)-f(y)|\ge\varepsilon_0.
  \]
\end{enumerate}

拼起来：

\[
\exists\varepsilon_0>0,\; \forall\delta>0,\;
\exists x,y\in D,\quad
|x-y|<\delta\;\land\;|f(x)-f(y)|\ge\varepsilon_0.
\]

翻译回日常语言：存在某个固定的距离 \(\varepsilon_0\)，使得无论你把 \(\delta\) 缩到多小，总能在定义域里找到两个点，它们靠得比 \(\delta\) 还近，但函数值却差了至少 \(\varepsilon_0\)。这是``不一致连续''的精确定义。等后面学到实数分析的时候，这个否定形式会和原定义一样重要。

否定复杂命题不要靠语感硬翻。中文里``不是没有''等于``有''，英语里双重否定可以约掉，但量词嵌套叠到三层以上，语感完全不可靠。老老实实从最外层一层一层往里推，每步只翻一层。慢，但不出错。

\subsection{量词没有告诉你的东西}

\begin{itemize}
\tightlist
\item
  \(\forall x\in D,P(x)\) 为真，不等于你能\textbf{实际检查}每一个元素。对无限集合，全称命题的证明靠的是论证结构，不靠逐个枚举。
\item
  \(\exists x\in D,P(x)\) 为真，不等于你知道\textbf{具体是哪一个}。很多存在性证明只论证``不可能一个都没有''，却不给出构造——这叫非构造性存在证明，在数学里合法，但在编程里不够用。
\item
  量词告诉你的只是``有没有''和``有多少个满足''，它不告诉你这些元素的\textbf{分布结构}。比如 \(\forall x\exists y\) 为真，不说明 \(y\) 和 \(x\) 之间是连续的还是跳跃的、是线性的还是混沌的。
\item
  论域的选择本身就是一个建模决定，量词形式化不会替你审查这个决定对不对。把论域设成``所有人''还是``所有收到问卷的人''，结论可能完全相反，而量词只管封装好的论域里的事。
\end{itemize}

后面到了实数分析和概率论，嵌套量词会密集出现。那时候再回来看这一节的翻转规则，它就是你的安全带。

\subsection{参考来源}

\begin{itemize}
\tightlist
\item
  MIT Mathematics for Computer Science 的 Predicates and Quantifiers 相关章节。
\item
  Discrete Mathematics: An Open Introduction 的 Logic and Proofs 部分。
\end{itemize}
